A foundation model-based agent operates in a perception--reasoning--action loop, using an LLM to process observations and execute actions such as tool invocations, code execution, and inter-agent communication. The underlying model workloads can also be multimodal: recent systems combine visual, audio, and clinical or platform metadata for domain-specific prediction and governance~\cite{wang-etal-2025-reasoning-enhanced,10.1145/3711896.3737195,xu2026mobifusiondrm}. A multi-agent system (MAS) comprises multiple such agents, as exemplified by MetaGPT~\cite{metagpt2024}, CAMEL~\cite{camel2023}, ChatDev~\cite{chatdev2024}, and IoA~\cite{ioa2025}, operating within a shared environment or communicating to accomplish collective objectives~\cite{li2024, wang2026classical}. Recent multi-agent applications also span reinforcement-learning-enhanced collaborative reasoning and multimodal misinformation detection~\cite{li2026draftrl,li2026retrievalaugmentedfakenews}. MAST~\cite{mast2025} identifies 14 recurring failure patterns in such systems, motivating runtime diagnostics rather than evaluation only at the final task output.

Representation-level adaptation is another adjacent concern: recent speech work studies adversarial attribute obfuscation, accent-controlled editing, and pronunciation representations~\cite{11249000,qu2026partialaccentcontroleditingfrozen,qu2026phonemeawarepronunciationrepresentationsl2english}. These works are not multi-agent infrastructure methods, but they illustrate the heterogeneity of model state and representations that agentic applications may invoke.

Three broad \emph{system properties} recur across the literature reviewed here~\cite{understandingmas2026}: (1)~distributed execution across multiple agents; (2)~inter-agent communication, which creates traces spanning agent boundaries and allows failures to propagate through communication channels~\cite{famas2025, traceelephant2026}; and (3)~inter-agent memory sharing, where prompts, tools, and cached state may be reused across the system. Section~\ref{sec:infra_challenges} derives a serving-side view from these properties and highlights three workload characteristics---sparse activation, invocation distance, and inter-agent memory sharing---that directly affect memory placement, scheduling, and cache policy. This distinction separates general MAS structure from the narrower workload characteristics used in our infrastructure analysis.

\paragraph*{Review scope, selection, and evidence coding.} This article is a curated critical review, not a systematic review. We focus on runtime operational mechanisms for foundation-model agents and MAS, distinct from three well-surveyed areas: orchestration~\cite{zhu2026orchestration}, single-model inference~\cite{kwon2023}, and offline evaluation~\cite{evalbenchmark2025}. For inclusion in the system tables, a work must directly address at least one runtime function in our scope: failure localization, observability/tracing, safety intervention, diagnostic benchmarking, or serving-time memory, scheduling, cache, or power management for agentic workloads. General LLM-serving systems are included when they provide a baseline or mechanism directly reused by agent workloads; orchestration-only frameworks, general agent architectures, and long-term memory representation without serving-time resource management are used only as context. To make the review more reproducible without overstating exhaustiveness, candidate papers were collected iteratively from primary systems papers, adjacent surveys, major peer-reviewed venues, and forward/backward citation tracing. We prioritize work from 2024--2026 because foundation-model multi-agent runtime research is unusually recent, while retaining older serving papers when they define mechanisms directly reused by agentic systems. A paper enters the core system tables only when its runtime role can be identified from the source and mapped to one of the functions above. Because the corpus was not generated by a preregistered database protocol with fixed search strings and PRISMA-style screening, the reported counts should be interpreted as a curated landscape rather than an exhaustive census.

To make the evidential depth of individual claims explicit, we assign cited systems to the following tiers:

\begin{enumerate}
\item \textbf{Tier 1a (full-text verified).} The source was read in full and the specific quantitative claim was located in the body text or a table. These numbers may be compared and quoted when the underlying benchmark and protocol support the comparison.
\item \textbf{Tier 1b (read in full, not page-verified).} The source was read in full and has a confirmed peer-reviewed venue, but the quantitative claim was taken from the published text without being cross-checked against a specific page, table, or figure. These numbers may be quoted with the $^\dagger$ marker but are not used to compute cross-system deltas.
\item \textbf{Tier 2 (abstract-only).} Only the abstract was read; the body was not examined. This applies equally to preprints and venue-published papers whose body we did not read. The claim is reported as stated and marked \emph{[a]}. The distinction from Tier~1b is verification depth, not venue status.
\item \textbf{Tier 3.} Documentation, platform pages, and vendor reports, labelled by source type and used for context only.
\end{enumerate}

Applying this procedure yields \textbf{8 Tier-1 papers across 9 table entries: 5 in Tier~1a} (FAMAS, TrajAudit, Agent Traces, TraceElephant, ScaleSim) and \textbf{3 in Tier~1b} (G-Safeguard at ACL 2025, ShieldAgent at ICML 2025, and AgenTracer at ICLR 2026). RootSE is the benchmark introduced by TrajAudit, so the two rows share one paper. Systems such as RAFFLES, AgentLocate, MultiAgentBench, and Who\&When have confirmed venues but were read at abstract level and therefore remain Tier~2; LlamaFirewall is likewise Tier~2. Throughout the review, publication venue and verification tier are kept separate, and quantitative values from heterogeneous benchmarks are presented as landscape evidence rather than as a global ranking.
